\chapter{From repetition to regularity}
\chaptertag{Laws of large numbers, the central limit theorem, and sampling}

\begin{question}
Why do averages stabilize, and why do bell curves emerge from non-normal data?
\end{question}

Let $X_1,\ldots,X_n$ be independent copies of a variable with mean $\mu$ and variance $\sigma^2$. Their average is
\[
\bar X_n=\frac1n\sum_{i=1}^nX_i,
\qquad
\E[\bar X_n]=\mu,
\qquad
\Var(\bar X_n)=\frac{\sigma^2}{n}.
\]

\begin{theorem}[Weak law of large numbers]
For every $\varepsilon>0$,
\[
\P(|\bar X_n-\mu|>\varepsilon)\longrightarrow0.
\]
\end{theorem}

With finite variance, Chebyshev gives the proof immediately:
\[
\P(|\bar X_n-\mu|>\varepsilon)
\le \frac{\sigma^2}{n\varepsilon^2}.
\]

\begin{theorem}[Central limit theorem]
If $0<\sigma^2<\infty$, then
\[
\frac{\sqrt n(\bar X_n-\mu)}{\sigma}
\xrightarrow{d}\Normal(0,1).
\]
\end{theorem}

The LLN says where the average goes; the CLT describes the scale and approximate shape of its remaining fluctuations.

\begin{center}
\begin{tikzpicture}[x=.7cm,y=.8cm]
  \foreach \x/\h in {0/1.0,1/3.4,2/1.2,3/3.1,4/.8}{\fill[rose!65] (\x,0) rectangle +(0.7,\h);}
  \node at (1.9,-.6) {one draw};
  \draw[-{Stealth},thick,mint] (4.3,1.8)--(5.5,1.8);
  \foreach \x/\h in {6/.4,7/1.2,8/2.3,9/3.1,10/2.3,11/1.2,12/.4}{\fill[violet!65] (\x,0) rectangle +(0.7,\h);}
  \node at (9.3,-.6) {average of many draws};
\end{tikzpicture}
\end{center}

\section{Sampling distributions}

A \emph{statistic} is a function of the sample. Before data are observed it is random, and its probability law is its sampling distribution. The standard deviation of that distribution is the \emph{standard error}.

For a Bernoulli sample, the sample proportion $\hat p$ has
\[
\E[\hat p]=p,\qquad
\operatorname{SE}(\hat p)=\sqrt{\frac{p(1-p)}n}.
\]

\section{Monte Carlo and bootstrap}

Monte Carlo approximates an expectation by simulating from a known model. The bootstrap approximates a sampling distribution by repeatedly resampling, with replacement, from observed data. The bootstrap is powerful but cannot repair a badly biased or unrepresentative original sample.

\begin{exercisebox}
\begin{enumerate}
  \item Compute the standard error of the mean when $\sigma=12$ and $n=64$. \exlevel{Fluency}
  \item Prove the weak law using Chebyshev's inequality. \exlevel{Reasoning}
  \item Approximate $\P(480\le X\le520)$ for $X\sim\Bin(1000,.5)$ using the CLT and a continuity correction. \exlevel{Fluency}
  \item Explain why quadrupling the sample size halves a standard error. \exlevel{Reasoning}
  \item Simulate sample means from a skewed distribution for $n=1,2,5,30$. Quantify convergence rather than only describing the pictures. \exlevel{Exploration}
  \item Compare Monte Carlo error at $n=100$, $10^4$, and $10^6$. Is the $n^{-1/2}$ rate visible? \exlevel{Exploration}
  \item Give an example where bootstrap resampling cannot reveal a sampling bias. \exlevel{Reasoning}
\end{enumerate}
\end{exercisebox}
